%-------------------------
% Resume in Latex
% Author : Jake Gutierrez
% Copyright (c) 2020 Jake Gutierrez
% Based off of: https://github.com/sb2nov/resume
% License : MIT -- full text in NOTICE
%------------------------

% Dense single-column layout. Implements the command interface in
% references/INTERFACE.md. Input this from a driver, after the driver has set
% its identity macros and before \begin{document}.

\usepackage{latexsym}
\usepackage[empty]{fullpage}
\usepackage{titlesec}
\usepackage{marvosym}
\usepackage[usenames,dvipsnames]{color}
\usepackage{verbatim}
\usepackage{enumitem}
\usepackage[hidelinks]{hyperref}
\usepackage{fancyhdr}
\usepackage[english]{babel}
\usepackage{tabularx}
\usepackage{fontawesome5}
\usepackage{multicol}
\setlength{\multicolsep}{-3.0pt}
\setlength{\columnsep}{-1pt}
\input{glyphtounicode}

% Serif or sans. The driver must declare this before \input-ing the template:
%   \newif\ifResumeTimes
%   \ResumeTimesfalse      % or \ResumeTimestrue for serif
% It cannot be defaulted here: a \newif guarded by \ifdefined breaks TeX's
% \if/\fi matching when the conditional already exists.
\ifResumeTimes
  \usepackage{newtxtext}
\else
  \usepackage[default]{sourcesanspro}
\fi

% Spacing knobs. A driver may override these before \input-ing this file; two
% roles are often tuned to slightly different values.
\providecommand{\ResumeSectionSkip}{-7pt}
\providecommand{\ResumeProjectTitleSkip}{3pt}

\pagestyle{fancy}
\fancyhf{}
\fancyfoot{}
\renewcommand{\headrulewidth}{0pt}
\renewcommand{\footrulewidth}{0pt}

\addtolength{\oddsidemargin}{-0.6in}
\addtolength{\evensidemargin}{-0.5in}
\addtolength{\textwidth}{1.19in}
\addtolength{\topmargin}{-.7in}
\addtolength{\textheight}{1.4in}

\urlstyle{same}
\raggedbottom
\raggedright
\setlength{\tabcolsep}{0in}

\titleformat{\section}{
  \vspace{\ResumeSectionSkip}\raggedright\Large\bfseries
}{}{0em}{}[\color{black}\titlerule \vspace{-2pt}]

% Machine readable / ATS parsable output
\pdfgentounicode=1

%-------------------------
% Command interface
\newcommand{\resumeItem}[1]{
  \item\small{
    {#1 \vspace{-2pt}}
  }
}

\newcommand{\resumeSubheading}[4]{
  \vspace{-2pt}\item
    \begin{tabular*}{1.0\textwidth}[t]{l@{\extracolsep{\fill}}r}
      \textbf{\normalsize #1} & \textbf{\small #2} \\
      \textit{\small#3} & \textit{\small #4} \\
    \end{tabular*}\vspace{-7pt}
}

% This layout renders all three the same way. Others do not -- keep them
% distinct so a new template can differentiate without forking content.
\newcommand{\resumeEducation}[3]{\resumeSubheading{#1}{#2}{#3}{}}
\newcommand{\resumePublication}[3]{\resumeSubheading{#1}{#2}{#3}{}}

\newcommand{\resumeSubSubheading}[2]{
    \item
    \begin{tabular*}{0.97\textwidth}{l@{\extracolsep{\fill}}r}
      \textit{\small#1} & \textit{\small #2} \\
    \end{tabular*}\vspace{-7pt}
}

\newcommand{\resumeProjectHeading}[2]{
    \item
    \begin{tabular*}{1.001\textwidth}{l@{\extracolsep{\fill}}r}
      \textbf{\normalsize #1} & \textbf{\small #2}\\
    \end{tabular*}\vspace{-7pt}
}

\newcommand{\resumeSubItem}[1]{\resumeItem{#1}\vspace{-4pt}}

\newcommand{\resumeProjectTitle}[1]{
  \item[]
  \vspace{\ResumeProjectTitleSkip}
  \small\textbf{#1}
  \vspace{-1pt}
}

\renewcommand\labelitemi{$\vcenter{\hbox{\tiny$\bullet$}}$}
\renewcommand\labelitemii{$\vcenter{\hbox{\tiny$\bullet$}}$}

\newcommand{\resumeSubHeadingListStart}{\begin{itemize}[leftmargin=0.0in, label={}]}
\newcommand{\resumeSubHeadingListEnd}{\end{itemize}}
\newcommand{\resumeItemListStart}{\begin{itemize}}
\newcommand{\resumeItemListEnd}{\end{itemize}\vspace{-5pt}}
\newcommand{\resumeItemListStartTight}{\begin{itemize}[topsep=0pt]}

\newcommand{\resumeSectionGap}{\vspace{-16pt}}

% This layout's metrics are tight enough to want hand-tuned negative leading.
\newcommand{\resumeTighten}[1]{\vspace{#1}}

% Identity defaults. Drivers set these via profiles/<person>.tex.
\providecommand{\ResumeName}{Your Name}
\providecommand{\ResumeEmail}{you@example.com}
\providecommand{\ResumePhone}{000-000-0000}
\providecommand{\ResumeLinkedInName}{Your Name}
\providecommand{\ResumeLinkedInURL}{https://www.linkedin.com/in/your-handle/}

%----------HEADING----------
\newcommand{\ResumeHeader}{%
\begin{center}
    {\Huge  \ResumeName} \\  \vspace{1pt}


     \small \href{mailto:\ResumeEmail}{\raisebox{-0.2\height}\faEnvelope\  \underline{\ResumeEmail}} ~
     ~
 {\raisebox{-0.2\height}\faPhone\ {\ResumePhone}}
     ~
    \href{\ResumeLinkedInURL}{\raisebox{-0.2\height}\faLinkedin\ \underline{\ResumeLinkedInName}}


    \vspace{-8pt}
\end{center}
}
